%%%%%%%%%%  scan idx 187  |  folio 174  %%%%%%%%%%
est
pleinement fidèle et a pour image essentielle $\mathrm D(\Qcoh(S))_{\parf}$ (2.2.1).
\end{enumerate}

\emph{Preuve.} Soit $F\to G$ un épimorphisme de $\Qcoh(S)$, où $G$ est de type fini. Comme $F$ est limite inductive de ses sous-faisceaux quasi-cohérents de type fini (EGA I 9.4.9), il existe, en vertu de (2.1.1 a)) un sous-faisceau quasi-cohérent de type fini $F'$ de $F$ tel que le composé $F'\to F\to G$ soit un épimorphisme. D'après (2.2.3 (ii)), il existe un épimorphisme $E\to F'$, avec $E\in\Ob L(S)$, donc le composé $E\to F'\to F$, composé avec $F\to G$ est un épimorphisme, ce qui prouve que la catégorie des $\mathcal O_S$-Modules quasi-cohérents de type fini est quasi-relevable dans $\Qcoh(S)$, donc a fortiori il en est de même de la catégorie $\Loclib(S)$ ou $L(S)$. Ceci démontre la première assertion de a). Pour la seconde, appliquons (1.1) en prenant pour $S$ le site zariskien de $S$, pour $\mathcal C$ (resp. $\mathcal C_0$) la catégorie fibrée $U\mapsto\Qcoh(U)$ (resp. $L(U)$) ($U$ désigne un ouvert variable de $S$ et $L(U)$ la sous-catégorie additive strictement pleine de $\Qcoh(U)$ engendrée par les $(L_i^{\otimes-n}\,|\,U)$). Comme tout objet de $L$ est pseudo-cohérent, les notions de $n$-pseudo-cohérence (resp. pseudo-cohérence) (2.2.1) dans la catégorie fibrée $\mathrm D(\Qcoh)$ coïncident avec les mêmes notions relativement à $L$ (I 2.14). La condition (ii) (non respée) de (1.1) étant satisfaite en vertu de (2.2.3 (ii)), on en déduit, grâce à (1.1 a)), la seconde partie de a) relative à $L(S)$ ; l'assertion relative à $\Loclib(S)$ se prouve de manière analogue. Les assertions b) découlent immédiatement de (1.1 b) et c)), compte tenu du fait que si $E\in\Ob\mathrm D(\Qcoh(S))$ est $n$-pseudo-cohérent, alors $E$ est objet de $\mathrm D^{-}(\Qcoh(S))$ puisque $S$ est quasi-compact. Un objet de la catégorie fibrée $\Qcoh$ est localement dans $L$ si et seulement s'il est localement libre de type fini : la catégorie $L$ est donc localement relevable dans $\Qcoh$ et localement stable par noyau d'épimorphisme. Par suite, les assertions c)) résultent de (1.2) appliqué pour $\mathcal C=\Qcoh$, $\mathcal C_0=L$, $\mathcal C_{1S}=\Loclib(S)$. Ceci achève la démonstration

\textbf{Proposition 2.2.9.} Les hypothèses et notations sont celles de (2.2.8), mais on suppose en outre $S$ séparé ou noethérien.
\begin{enumerate}[label=\alph*)]
\item Les foncteur canoniques
\[
\mathrm K^{-,b}(L(S))/\mathrm K^{-,\varnothing}(L(S))\longrightarrow\mathrm D(S)
\]
et
\[
\mathrm K^{-,b}(\Loclib(S))/\mathrm K^{-,\varnothing}(\Loclib(S))\longrightarrow\mathrm D(S)
\]
